% ============================================================================
% sec7_counterfactual.tex — Paper 2: Welfare counterfactuals
% ============================================================================

\section{Welfare Counterfactuals}
\label{sec:counterfactual}

The calibrated model (Section~\ref{sec:results_calibration}) allows
us to compute counterfactuals that flow directly from the structural
parameters, rather than relying solely on reduced-form price
associations. The counterfactual chain is:
\emph{policy change} $\to$ \emph{change in $m^*$} (via
Proposition~\ref{prop:optimal_m} and calibrated primitives) $\to$
\emph{change in FL incidence} $\to$ \emph{change in procurement
spending} (via reduced-form price association). The last link
uses conditional correlations, not causal effects; to the extent the
price associations reflect selection, the welfare gains are
overstated. We present them as order-of-magnitude benchmarks.

\paragraph{CF1: Removing the minimum-bidder requirement.}
Under convite, 39.2\% of FL-present tenders have
$n_{\text{genuine}} < 3$, forcing cover-bidder deployment (corner
solution, Proposition~\ref{prop:optimal_m}); 23.5\% have
$n_{\text{genuine}} = 0$ (pure cover bidding---every participant is
FL). The calibrated interior optimum for these tenders is
$\hat{m}_{\text{interior}} = 1.4$, well below the observed corner
mean of 2.8 (Section~\ref{sec:results_calibration}). Removing the
minimum-bidder rule replaces the corner solution with the interior
optimum: for tenders where $\hat{m}_{\text{interior}} = 0$ (i.e.,
cover bidding is unprofitable absent the constraint), FL deployment
ceases entirely. The calibrated model predicts a 62\% reduction in
FL incidence within convite, yielding illustrative welfare gains of
R\$74--211M over 2009--2019.%
\footnote{The range reflects the reduced-form price coefficient used
in the last link: R\$74M uses the cross-fit lower bound (0.036);
R\$211M uses the IV upper bound (0.194).}

This counterfactual has a concrete institutional pathway:
Brazil's Lei~14.133/2021 (the ``New Procurement Law'') relaxes the
minimum-bidder constraint for preg\~ao but retains it for convite.
Extending the relaxation to convite would eliminate the institutional
demand for cover bidding at zero administrative cost.

\paragraph{CF2: Welfare-maximizing detection threshold.}
We embed the ROC curve in a social welfare function:
\begin{equation}
  W(t) = \delta \cdot \hat{\beta} \cdot S_{\text{FL}}(t) -
  c_{\text{inv}} \cdot N_{\text{flagged}}(t)
  \label{eq:welfare_function}
\end{equation}
where $t$ is the IQR multiplier, $\delta$ is the deterrence rate,
$\hat{\beta}$ is the price association, $S_{\text{FL}}(t)$ is total
spending in FL-flagged tenders at threshold $t$, and $c_{\text{inv}}$
is the per-firm investigation cost. The deterrence rate $\delta$
enters through the structural model: detection raises $\theta_k$,
which reduces $m^*$ via the calibrated $\hat{\phi}_0 / \hat{c}_1$
ratio.

At R\$100K/investigation with $\delta = 0.50$, the optimal
threshold is $t^* = 1.6\times$IQR (net welfare R\$135M; 2,537
flagged firms)---close to the baseline $1.5\times$ and the
Youden-optimal $1.45\times$
(Figure~\ref{fig:counterfactual_screening}). The convergence of
three independent threshold criteria---statistical (IQR rule),
detection-theoretic (Youden's $J$), and welfare-maximizing
($W(t)$)---at approximately $1.5\times$IQR provides a coherent
basis for operational deployment.

Screening becomes welfare-negative at $c_{\text{inv}} >
\text{R\$500K}$. For competition authorities with low-cost desk
review capacity (R\$50--100K per firm), the FL screen generates
substantial positive welfare at every threshold between
$1.0\times$ and $3.0\times$IQR.

\begin{figure}[htbp]
\centering
\includegraphics[width=0.82\textwidth]{images/fig_counterfactual_screening.pdf}
\caption{Welfare-maximizing threshold vs.\ investigation cost, by
deterrence rate $\delta$. At R\$100K, $t^* \approx 1.6\times$IQR.}
\label{fig:counterfactual_screening}
\end{figure}

\paragraph{CF3: Increasing detection probability.}
The model predicts $\partial m^* / \partial \theta_k < 0$: higher
detection probability deters cover-bidder deployment. Empirically,
mean FL per tender decreases from Q1 (0.41) to Q4 (0.31) by PBU
size, consistent with the sign prediction. However, the calibrated
detection-penalty parameter $\hat{\phi}_0 = 0.019$ is
indistinguishable from zero (CI~$[0.000, 0.056]$;
Table~\ref{tab:cartel_primitives}), limiting the precision of
quantitative counterfactuals along this margin. Using the point
estimate as an upper bound, a 10\% increase in detection probability
would reduce $m^*$ by at most 0.04 firms per tender---a modest
effect that reinforces the calibration finding: cartel cover-bidding
intensity responds primarily to competitive pressure ($\hat{\gamma}$),
not to enforcement ($\hat{\phi}_0$). The policy implication is that
detection-based screening (CF2) dominates deterrence-based
investment in oversight capacity.

../output/tables/tab_counterfactual_welfare.tex
